\documentclass[11pt,a4paper]{article}
\usepackage[utf8]{inputenc}
\usepackage[T1]{fontenc}
\usepackage[english]{babel}
\usepackage{graphicx}
\usepackage{amsmath}
\usepackage{amssymb}
\usepackage{booktabs}
\usepackage{hyperref}
\usepackage{listings}
\usepackage{xcolor}
\usepackage{float}
\usepackage{geometry}
\usepackage{fancyhdr}
\usepackage{caption}
\usepackage{subcaption}
\usepackage{url}

\geometry{margin=2.5cm}

% Code listing style
\definecolor{codebg}{rgb}{0.95,0.95,0.95}
\definecolor{codegreen}{rgb}{0,0.5,0}
\definecolor{codegray}{rgb}{0.5,0.5,0.5}
\definecolor{codepurple}{rgb}{0.5,0,0.33}
\definecolor{codeblue}{rgb}{0,0,0.7}

\lstdefinestyle{mystyle}{
    backgroundcolor=\color{codebg},
    commentstyle=\color{codegreen},
    keywordstyle=\color{codeblue},
    numberstyle=\tiny\color{codegray},
    stringstyle=\color{codepurple},
    basicstyle=\ttfamily\footnotesize,
    breakatwhitespace=false,
    breaklines=true,
    captionpos=b,
    keepspaces=true,
    numbers=left,
    numbersep=5pt,
    showspaces=false,
    showstringspaces=false,
    showtabs=false,
    tabsize=2,
    frame=single,
    framesep=3pt,
    float=H,
}
\lstset{style=mystyle}

\title{\textbf{FarmifAI: An Offline Agricultural AI Assistant for Rural Communities}}
\author{}
\date{December 2025}

\begin{document}

\maketitle

\begin{abstract}
This paper presents FarmifAI, a fully offline mobile application designed to provide agricultural assistance to farmers in rural areas with limited or no internet connectivity. The system integrates four AI models: (1) a vision model based on MobileNetV2 for plant disease classification across 21 disease classes in 5 crops, (2) an NLP model using MiniLM for semantic search with 384-dimensional embeddings, (3) a quantized Llama 3.2 1B language model for contextual response generation, and (4) Vosk-based speech recognition for hands-free voice interaction. All inference is performed on-device using MindSpore Lite and llama.cpp, achieving response times under 5 seconds on mid-range Android devices. The application uses semantic search with a knowledge base to provide accurate, context-aware responses to agricultural queries. Field testing demonstrates 92.3\% accuracy in disease detection and high user satisfaction among Colombian coffee farmers.
\end{abstract}

\textbf{Keywords:} Mobile AI, Offline Inference, Agricultural Technology, Plant Disease Detection, Semantic Search, Edge Computing

\section{Introduction}

Small-scale farmers in developing regions often lack access to reliable internet connectivity and expert agricultural advice. This technological gap can lead to significant crop losses due to undetected diseases and suboptimal farming practices. FarmifAI addresses this challenge by deploying a comprehensive AI assistant that operates entirely offline on Android devices.

The main contributions of this work are:
\begin{itemize}
    \item A multi-modal AI system combining vision, NLP, and voice capabilities for agricultural assistance
    \item An efficient on-device inference pipeline achieving real-time performance on consumer hardware
    \item A two-phase transfer learning approach for plant disease classification with limited training data
    \item Integration of semantic search with local LLM for accurate domain-specific responses
\end{itemize}

\section{System Architecture}

\subsection{Overview}

FarmifAI consists of four main components integrated into a single Android application (Figure~\ref{fig:system-arch}):

\begin{enumerate}
    \item \textbf{Voice Interface}: Vosk-based speech-to-text (STT) and Android Text-to-Speech (TTS) for hands-free operation
    \item \textbf{Vision Module}: MobileNetV2-based classifier for plant disease detection
    \item \textbf{Semantic Search}: MiniLM encoder with pre-computed knowledge base embeddings
    \item \textbf{Language Model}: Quantized Llama 3.2 1B for response generation
\end{enumerate}

\begin{figure}[H]
    \centering
    \includegraphics[width=0.85\textwidth]{fig_system_architecture.pdf}
    \caption{FarmifAI system architecture showing on-device components and data flow}
    \label{fig:system-arch}
\end{figure}

\subsection{Technology Stack}

Table~\ref{tab:tech-stack} summarizes the technologies used in each component.

\begin{table}[H]
    \centering
    \caption{Technology stack for FarmifAI components}
    \label{tab:tech-stack}
    \begin{tabular}{@{}llll@{}}
        \toprule
        \textbf{Component} & \textbf{Technology} & \textbf{Model} & \textbf{Size} \\
        \midrule
        Speech-to-Text & Vosk & model-es-small & $\sim$50MB \\
        Text-to-Speech & Android TTS & System engine & Built-in \\
        Semantic Search & MindSpore Lite & MiniLM (384-dim) & $\sim$90MB \\
        Plant Disease & MindSpore Lite & MobileNetV2 & $\sim$14MB \\
        Language Model & llama.cpp & Llama 3.2 1B Q4\_K\_M & $\sim$750MB \\
        \bottomrule
    \end{tabular}
\end{table}

\section{Vision Model: Plant Disease Classification}

\subsection{Model Implementation}

The vision model uses MobileNetV2~\cite{sandler2018mobilenetv2} as the backbone architecture, chosen for its efficiency on mobile devices. The model classifies plant images into 21 disease classes across 5 crops: coffee, corn, potato, pepper, and tomato.

\begin{lstlisting}[language=Python, caption={MobileNetV2 architecture with custom classification head}]
from tensorflow.keras.applications import MobileNetV2
from tensorflow.keras.layers import Dense, GlobalAveragePooling2D
from tensorflow.keras.layers import Dropout, BatchNormalization
from tensorflow.keras.models import Model
from tensorflow.keras.regularizers import l2

# Base model with ImageNet weights
base_model = MobileNetV2(
    input_shape=(224, 224, 3),
    include_top=False,
    weights='imagenet'
)

# Custom classification head
x = GlobalAveragePooling2D()(base_model.output)
x = BatchNormalization()(x)
x = Dropout(0.4)(x)
x = Dense(512, activation='relu', kernel_regularizer=l2(0.01))(x)
x = Dropout(0.4)(x)
output = Dense(21, activation='softmax')(x)

model = Model(inputs=base_model.input, outputs=output)
\end{lstlisting}

\subsection{Training Pipeline}

We employ a two-phase transfer learning strategy (Figure~\ref{fig:training-arch}) to maximize performance with limited agricultural data:

\nopagebreak
\textbf{Phase 1: Feature Extraction (25 epochs)}
\begin{itemize}
    \item Freeze all MobileNetV2 backbone layers
    \item Train only the classification head
    \item Learning rate: $10^{-3}$ with Adam optimizer
    \item Dropout rate: 0.4
\end{itemize}

\nopagebreak
\textbf{Phase 2: Fine-tuning (35 epochs)}
\begin{itemize}
    \item Unfreeze all layers for end-to-end training
    \item Learning rate: $10^{-5}$ (100x reduction)
    \item Label smoothing: 0.1
    \item Early stopping with patience of 7 epochs
\end{itemize}

\begin{figure}[H]
    \centering
    \includegraphics[width=0.9\textwidth]{fig_two_phase_training.pdf}
    \caption{Two-phase transfer learning strategy for MobileNetV2}
    \label{fig:two-phase}
\end{figure}

\begin{lstlisting}[language=Python, caption={Data augmentation configuration}]
train_datagen = ImageDataGenerator(
    rescale=1./255,
    rotation_range=40,
    width_shift_range=0.3,
    height_shift_range=0.3,
    shear_range=0.2,
    zoom_range=0.3,
    horizontal_flip=True,
    vertical_flip=True,
    brightness_range=[0.6, 1.4],
    fill_mode='reflect',
    validation_split=0.15
)
\end{lstlisting}

\begin{figure}[H]
    \centering
    \includegraphics[width=0.95\textwidth]{fig_training_architecture.pdf}
    \caption{Complete training architecture for vision and NLP models}
    \label{fig:training-arch}
\end{figure}

\subsection{Inference Pipeline}

On-device inference follows the pipeline shown in Figure~\ref{fig:vision-inference}:

\begin{figure}[H]
    \centering
    \includegraphics[width=0.95\textwidth]{fig_vision_inference.pdf}
    \caption{Vision model inference pipeline on Android device}
    \label{fig:vision-inference}
\end{figure}

\nopagebreak
\begin{lstlisting}[language=Kotlin, caption={Image preprocessing for MindSpore Lite inference}]
fun preprocessImage(bitmap: Bitmap): ByteBuffer {
    val scaled = Bitmap.createScaledBitmap(bitmap, 224, 224, true)
    
    // Allocate buffer: 224 * 224 * 3 channels * 4 bytes
    val buffer = ByteBuffer.allocateDirect(224 * 224 * 3 * 4)
    buffer.order(ByteOrder.nativeOrder())
    
    for (y in 0 until 224) {
        for (x in 0 until 224) {
            val pixel = scaled.getPixel(x, y)
            
            // Normalize to [-1, 1] range
            buffer.putFloat(((pixel shr 16 and 0xFF) - 127.5f) / 127.5f)
            buffer.putFloat(((pixel shr 8 and 0xFF) - 127.5f) / 127.5f)
            buffer.putFloat(((pixel and 0xFF) - 127.5f) / 127.5f)
        }
    }
    return buffer
}
\end{lstlisting}

\subsection{Deployment}

Model conversion from TensorFlow to MindSpore Lite format:

\begin{lstlisting}[language=bash, caption={Model conversion and deployment commands}]
# Export TensorFlow SavedModel
python tools/export_trained_model.py \
  --model_path outputs/best_model.h5 \
  --output_dir outputs/exported/

# Convert to MindSpore Lite format
python -m mindspore_lite.converter \
  --fmk=TF \
  --modelFile=outputs/exported/saved_model.pb \
  --outputFile=plant_disease_model \
  --inputShape=1,224,224,3

# Deploy to app assets
cp plant_disease_model.ms app/src/main/assets/
\end{lstlisting}

\section{NLP Model: Semantic Search}

\subsection{Model Implementation}

The semantic search component uses the \texttt{paraphrase-multilingual-MiniLM-L12-v2} model from Sentence Transformers to encode both queries and knowledge base entries into 384-dimensional embeddings.

\subsection{Embedding Generation}

\nopagebreak
\begin{lstlisting}[language=Python, caption={Mean pooling for sentence embeddings}]
import torch
import torch.nn.functional as F
from transformers import AutoTokenizer, AutoModel

def encode_with_mean_pooling(model, tokenizer, texts):
    inputs = tokenizer(
        texts, padding="max_length", truncation=True,
        max_length=128, return_tensors="pt"
    )
    
    with torch.no_grad():
        outputs = model(**inputs)
        attention_mask = inputs['attention_mask']
        token_embeddings = outputs.last_hidden_state
        
        # Expand mask for broadcasting
        input_mask_expanded = attention_mask.unsqueeze(-1).expand(
            token_embeddings.size()
        ).float()
        
        # Mean pooling
        sum_embeddings = torch.sum(
            token_embeddings * input_mask_expanded, 1
        )
        sum_mask = torch.clamp(input_mask_expanded.sum(1), min=1e-9)
        embeddings = sum_embeddings / sum_mask
        
        # L2 normalization
        embeddings = F.normalize(embeddings, p=2, dim=1)
    
    return embeddings.numpy()
\end{lstlisting}

\subsection{Query Processing Pipeline}

\nopagebreak
The semantic search pipeline (Figure~\ref{fig:rag}) processes user queries as follows:

\begin{enumerate}
    \item Tokenize input query (max 128 tokens)
    \item Encode query using MiniLM encoder
    \item Compute cosine similarity against 517 pre-computed KB embeddings
    \item Retrieve top-3 contexts above threshold (0.4)
    \item Augment prompt with retrieved contexts
    \item Generate response using Llama 3.2
\end{enumerate}

\begin{figure}[H]
    \centering
    \includegraphics[width=0.95\textwidth]{fig_rag_pipeline.pdf}
    \caption{Semantic search pipeline for context-aware response generation}
    \label{fig:rag}
\end{figure}

\begin{lstlisting}[language=Kotlin, caption={Cosine similarity search implementation}]
fun findTopKContexts(query: String, topK: Int = 3): List<Context> {
    val queryEmbedding = computeEmbedding(query)
    
    return kbEmbeddings.mapIndexed { i, emb ->
        val score = cosineSimilarity(queryEmbedding, emb)
        i to score
    }.filter { it.second >= 0.4f }
     .sortedByDescending { it.second }
     .take(topK)
     .map { (idx, score) ->
         Context(
             answer = kbEntries[idx].answer,
             score = score,
             category = kbEntries[idx].category
         )
     }
}

fun cosineSimilarity(a: FloatArray, b: FloatArray): Float {
    // Vectors are L2-normalized, dot product = cosine similarity
    return a.indices.sumOf { (a[it] * b[it]).toDouble() }.toFloat()
}
\end{lstlisting}

\section{LLM: Local Language Model}

\subsection{Model Implementation}

FarmifAI uses Llama 3.2 1B Instruct with Q4\_K\_M quantization, reducing model size from $\sim$2GB to $\sim$750MB while maintaining response quality.

\subsection{Inference}

\begin{lstlisting}[language=Kotlin, caption={Llama 3.2 prompt construction and generation}]
fun buildPrompt(context: String, query: String): String {
    val systemPrompt = """You are FarmifAI, an expert 
        agricultural assistant. Respond in Spanish.
        Use the provided context to answer questions."""
    
    return """<|begin_of_text|><|start_header_id|>system<|end_header_id|>

$systemPrompt<|eot_id|><|start_header_id|>user<|end_header_id|>

Context: $context
Question: $query<|eot_id|><|start_header_id|>assistant<|end_header_id|>

"""
}

suspend fun generate(prompt: String, maxTokens: Int = 150): Flow<String> = flow {
    llama.setTemperature(0.7f)
    llama.setTopP(0.9f)
    llama.generate(prompt).collect { token ->
        if (!token.contains("<|eot_id|>")) emit(token)
    }
}
\end{lstlisting}

\subsection{Deployment}

\begin{lstlisting}[language=bash, caption={LLM deployment to Android device}]
# Download quantized model
wget https://huggingface.co/bartowski/Llama-3.2-1B-Instruct-GGUF/\
resolve/main/Llama-3.2-1B-Instruct-Q4_K_M.gguf

# Deploy to device external storage
adb shell mkdir -p /sdcard/Android/data/edu.unicauca.app.agrochat/files/
adb push Llama-3.2-1B-Instruct-Q4_K_M.gguf \
  /sdcard/Android/data/edu.unicauca.app.agrochat/files/
\end{lstlisting}

\section{Voice Interface}

\subsection{Speech-to-Text (Vosk)}

\begin{lstlisting}[language=Kotlin, caption={Vosk STT integration}]
fun initializeVosk() {
    val modelDir = File(context.cacheDir, "model-es-small")
    if (!modelDir.exists()) {
        unpackAssets("model-es-small", modelDir)
    }
    voskModel = Model(modelDir.absolutePath)
}

fun startListening() {
    val recognizer = Recognizer(voskModel, 16000.0f)
    speechService = SpeechService(recognizer, 16000.0f)
    speechService?.startListening(object : RecognitionListener {
        override fun onResult(hypothesis: String?) {
            hypothesis?.let {
                val json = JSONObject(it)
                val text = json.getString("text")
                if (text.isNotBlank()) onResult?.invoke(text)
            }
        }
    })
}
\end{lstlisting}

\section{Results and Evaluation}

\subsection{Performance Metrics}

Table~\ref{tab:performance} shows the performance metrics for each component.

\begin{table}[H]
    \centering
    \caption{Performance metrics on mid-range Android device (Snapdragon 6xx)}
    \label{tab:performance}
    \begin{tabular}{@{}lcccc@{}}
        \toprule
        \textbf{Component} & \textbf{Latency} & \textbf{Memory} & \textbf{Accuracy} \\
        \midrule
        Vision Model & 200-400ms & 50MB & 92.3\% Top-1 \\
        Semantic Search & 60-130ms & 100MB & 94.4\% Top-1 \\
        LLM Generation & 2-5s & 1.2GB & -- \\
        Voice STT & Real-time & 100MB & 85-92\% \\
        \bottomrule
    \end{tabular}
\end{table}

\subsection{Disease Detection Results}

The vision model achieves:
\begin{itemize}
    \item \textbf{Top-1 Accuracy}: 92.3\%
    \item \textbf{Top-3 Accuracy}: 98.1\%
    \item \textbf{Inference Time}: 200-400ms
\end{itemize}

\section{Deployment Workflow}

Figure~\ref{fig:deployment} illustrates the complete deployment workflow from model training to device installation.

\begin{figure}[H]
    \centering
    \includegraphics[width=0.95\textwidth]{fig_deployment_workflow.pdf}
    \caption{Model deployment workflow}
    \label{fig:deployment}
\end{figure}

\section{Datasets and Data Sources}

\subsection{Vision Model Training Data}

The plant disease classification model was trained using two complementary datasets:

\textbf{PlantVillage Dataset:} A publicly available dataset containing 54,309 images of healthy and diseased plant leaves across 14 crop species. We extracted images for our target crops (tomato, potato, pepper, corn) covering 18 disease classes plus healthy states. Available at: \url{https://github.com/spMohanty/PlantVillage-Dataset} and \url{https://www.kaggle.com/datasets/emmarex/plantdisease}.

\textbf{BRACOL Coffee Dataset:} The Brazilian Arabica Coffee Leaf dataset (BRACOL) containing 1,747 images of arabica coffee leaves affected by the main biotic stresses. The dataset includes four disease classes:
\begin{itemize}
    \item Leaf miner (Minador)
    \item Leaf rust (Roya)
    \item Brown leaf spot
    \item Cercospora leaf spot
\end{itemize}
Available at: \url{https://data.mendeley.com/datasets/yy2k5y8mxg/1} (Paper: \url{https://doi.org/10.1016/j.compag.2019.105162}).

Both datasets were preprocessed to 224$\times$224 pixels and augmented using rotations, flips, zooms, and brightness variations to increase model robustness.

\subsection{NLP Knowledge Base}

The semantic search component uses a curated agricultural knowledge base consisting of:
\begin{itemize}
    \item 517 question-answer pairs covering crop management, disease identification, pest control, and best practices
    \item Questions sourced from agricultural extension services, farmer forums, and expert consultations
    \item Content validated by agronomists from the University of Cauca
    \item Categories: Crop Management (35\%), Disease Control (30\%), Pest Management (20\%), Soil \& Fertilization (15\%)
\end{itemize}

Each entry includes multiple question variations to improve semantic matching, with embeddings pre-computed using the paraphrase-multilingual-MiniLM-L12-v2 model.

\subsection{Voice Recognition Model}

The offline speech recognition uses Vosk's pre-trained Spanish language model (\texttt{vosk-model-small-es-0.42}):
\begin{itemize}
    \item Training data: Common Voice Spanish corpus + proprietary datasets
    \item Model size: 50MB compressed
    \item Optimized for Latin American Spanish dialects
    \item Adapted vocabulary includes agricultural terminology
\end{itemize}
Available at: \url{https://alphacephei.com/vosk/models/vosk-model-small-es-0.42.zip} (Official models: \url{https://alphacephei.com/vosk/models}).

\subsection{Language Model}

The Llama 3.2 1B Instruct model was developed by Meta AI:
\begin{itemize}
    \item Pre-trained on diverse internet text and code
    \item Fine-tuned for instruction-following tasks
    \item Quantized to Q4\_K\_M format using llama.cpp
    \item No additional domain-specific fine-tuning applied
\end{itemize}
Available at: \url{https://huggingface.co/meta-llama/Llama-3.2-1B-Instruct} (Quantized: \url{https://huggingface.co/bartowski/Llama-3.2-1B-Instruct-GGUF}).

\end{document}
